\documentclass[10pt]{article}
\usepackage[margin=0.8in]{geometry}
\usepackage{amsmath,amsfonts,enumitem}
\setlength{\parskip}{3pt}
\setlist{nosep}
% Paper labels, so references resolve here.
\makeatletter
\newlabel{eq:10}{{10}{3}}\newlabel{eq:11}{{11}{3}}\newlabel{tab:field_summary}{{I}{2}}
\makeatother
\begin{document}

\noindent\textbf{Proposed Section VI-B, replacing ``Comparison with Other Primitives''.} Drafted 2026-10-30. This draft:
\begin{itemize}
\item supersedes \texttt{Queued\_Edits/VI-B\_comparison\_rewrite.tex};
\item deletes the vector-sum comparison and its figure;
\item answers the AE (matched fixed-radius baselines), Reviewer~1 Comment~5 and Reviewer~3.
\end{itemize}
Numbers come from \texttt{VF\_Robot/experiments/revision\_metrics/baselines.py} and await your review. Equation numbers follow the current paper.

\subsection*{B. Comparison with Matched Baselines}

To separate the estimate from the controller, the attraction law (\ref{eq:10}) and the orbital law (\ref{eq:11}) were run with the same gains, robot model, and formation, driven by three estimates of the critical point.
\begin{itemize}
\item \emph{Oracle} uses the true critical point and isolates the controller.
\item \emph{Magnitude} uses only the sensed speeds $m_i = \|\mathbf{v}(\mathbf{p}_i)\|$. A plane fit gives their slope $\mathbf{g}$ and their value $m_c$ at the centroid, and $\hat{\mathbf{p}}_m = \mathbf{p}_c - m_c\mathbf{g}/\|\mathbf{g}\|^2$ locates the zero of a speed that grows linearly with distance. It uses the gains of the proposed laws without separate tuning.
\end{itemize}
Table~\ref{tab:baselines} reports both, alongside the plane fit, in simulation over repeated trials.

On the affine fields, the plane fit matches the oracle in every metric, because its estimate is exact. The remaining errors belong to the controller: the friction ring of radius $v_{\min}/k$ in attraction, and the outward lag offset of the orbit. On the reconstructed fields, the plane fit's estimate errs by 0.05--0.06~m, and its orbit error at $r_d = 0.3$~m rises above the oracle's.

The magnitude estimate errs by 0.06~m on the affine fields and 0.21~m on the reconstructed ones. Near the critical point the sensed speeds form a cone, the fitted slope vanishes, and the estimate diverges.
\begin{itemize}
\item \emph{Attraction.} The cluster never comes to rest, and after 10~s it still moves at about 0.12~m/s. Its smaller final error is a snapshot of that motion, and it costs about four times the control effort.
\item \emph{Orbit.} Its smaller errors at 0.3 and 0.5~m on the affine fields come from an inward estimate bias that offsets the lag. The same bias holds the orbit 0.07~m inside the commanded 0.1~m.
\item \emph{Type.} The canonical sink, source, vortex, and saddle share the speed field $\|\mathbf{v}\| = r$, so only the plane fit identifies the type of the critical point (Table~\ref{tab:field_summary}).
\end{itemize}

[PENDING the new campaign.] In hardware, the oracle runs of Section~V repeat each attraction start and orbital radius with the measured critical point, and give the controller's share of the hardware error directly.

\begin{table}[h]
\centering
\footnotesize
\caption{Matched baselines in simulation: the same attraction law (\ref{eq:10}) and orbital law (\ref{eq:11}), gains, robot model, and formation, driven by three estimates of the critical point. 500 attraction trials of 10~s and 50 orbit trials of 60~s per field and setting. Success is ending within 0.05~m; $\mu_e$ is the mean final error; $T_{\mathrm{conv}}$ is the median time to stay within 0.05~m; $E_u$ (m$^2$/s) is the realized quadratic effort; orbit RMSE of $r^* - r_d$ excludes the first 10~s; Est.\ err.\ is the mean $\|\hat{\mathbf{p}}^* - \mathbf{p}^*\|$ during the orbits.}
\label{tab:baselines}
\setlength{\tabcolsep}{2.5pt}
\begin{tabular}{llcccc|ccc|c}
\hline
 & & \multicolumn{4}{c|}{\textbf{Attraction}} & \multicolumn{3}{c|}{\textbf{Orbit RMSE (m)}} & \textbf{Est.} \\
\textbf{Field} & \textbf{Estimate} & \textbf{Succ.\ (\%)} & \textbf{$\mu_e$ (m)} & \textbf{$T_{\mathrm{conv}}$ (s)} & \textbf{$E_u$} & $r_d$ = 0.1 & $r_d$ = 0.3 & $r_d$ = 0.5 & \textbf{err.\ (m)} \\
\hline
Linear (6) & Plane fit & 100 & 0.025 & 2.1 & 0.074 & 0.093 & 0.051 & 0.033 & 0.000 \\
 & Oracle & 100 & 0.025 & 2.1 & 0.074 & 0.093 & 0.051 & 0.033 & 0.000 \\
 & Magnitude & 100 & 0.012 & 1.3 & 0.306 & 0.068 & 0.016 & 0.020 & 0.060 \\
\hline
Vortex (reconstructed) & Plane fit & 94 & 0.031 & 2.3 & 0.069 & 0.096 & 0.092 & 0.045 & 0.051 \\
 & Oracle & 100 & 0.025 & 2.1 & 0.074 & 0.093 & 0.051 & 0.033 & 0.000 \\
 & Magnitude & 94 & 0.034 & 1.8 & 0.269 & 0.065 & 0.079 & 0.068 & 0.211 \\
\hline
Saddle (reconstructed) & Plane fit & 100 & 0.025 & 2.3 & 0.071 & 0.090 & 0.071 & 0.050 & 0.061 \\
 & Oracle & 100 & 0.025 & 2.1 & 0.074 & 0.093 & 0.051 & 0.033 & 0.000 \\
 & Magnitude & 100 & 0.018 & 1.5 & 0.270 & 0.066 & 0.061 & 0.056 & 0.214 \\
\hline
\end{tabular}
\end{table}

\section*{Changes elsewhere}
\begin{enumerate}
\item \textbf{Delete} everything from ``The modular architecture allows navigation controllers'' through the \verb|\end{figure}| of \texttt{fig:primitive\_comparison} (\texttt{orbit\_demonstration.png}). Nothing else references the figure. Citations [22] and [48] remain cited elsewhere, so no bibliography entry is orphaned.
\item \textbf{Introduction roadmap.} ``Section VI compares performance against alternative navigation strategies and discusses limitations and future work.'' $\to$ ``Section VI compares the proposed estimate with matched baselines and discusses limitations and future work.''
\item \textbf{Superseded queued text.} Do not apply \texttt{Queued\_Edits/VI-B\_comparison\_rewrite.tex}, or item~10 of \texttt{prose\_framing\_batch.tex} (the vector-to-scalar description). This section describes the magnitude baseline itself.
\item \textbf{Metric definition} (item~7 of \texttt{prose\_framing\_batch.tex}). ``Enters and stays within $\varepsilon$'' is met by a cluster that keeps moving, as the magnitude baseline shows. Either require the cluster to come to rest (speed below a threshold for the last second), or report the final speed with the final error. The proposed laws always park under friction, so their numbers do not change.
\item \textbf{Hardware oracle runs} come from block~C of the campaign prompt. The bracketed sentence stays out until those runs exist.
\end{enumerate}

\textbf{Response letter (vector-sum only; the vector-to-scalar orbit was never run).}
\begin{quote}
We agree that the vector-sum controller lacks radial feedback, so its drift said little about the proposed method. We removed that comparison and replaced Section VI-B with matched baselines: the same attraction and orbital laws driven by the true critical point and by a magnitude-only estimate, scored with the metrics of Section V-B over repeated trials.
\end{quote}

\textbf{For Section IV, not this section.} On the nonlinear fields of the hyperparameter study, both estimates lose the orbit at $r_d = 0.5$~m in all trials (gyre and curved saddle), while the oracle holds it. Attraction on the gyre reaches a neighboring critical point in about 42\% of trials for both estimates. These belong with the nonlinear benchmarks.

\end{document}
